pdfoutput=1
% Options for packages loaded elsewhere
\PassOptionsToPackage{unicode}{hyperref}
\PassOptionsToPackage{hyphens}{url}
\pdfoutput=1
\documentclass[
  12pt,
]{article}
\usepackage{xcolor}
\usepackage{amsmath,amssymb}
\setcounter{secnumdepth}{-\maxdimen} % remove section numbering
\usepackage{iftex}
\ifPDFTeX
  \usepackage[T1]{fontenc}
  \usepackage{textcomp} % provide euro and other symbols
\else % if luatex or xetex
  \usepackage{unicode-math} % this also loads fontspec
  \defaultfontfeatures{Scale=MatchLowercase}
  \defaultfontfeatures[\rmfamily]{Ligatures=TeX,Scale=1}
\fi
\usepackage{lmodern}
\ifPDFTeX\else
  % xetex/luatex font selection
\fi
% Use upquote if available, for straight quotes in verbatim environments
\IfFileExists{upquote.sty}{\usepackage{upquote}}{}
\IfFileExists{microtype.sty}{% use microtype if available
  \usepackage[]{microtype}
  \UseMicrotypeSet[protrusion]{basicmath} % disable protrusion for tt fonts
}{}
\makeatletter
\@ifundefined{KOMAClassName}{% if non-KOMA class
  \IfFileExists{parskip.sty}{%
    \usepackage{parskip}
  }{% else
    \setlength{\parindent}{0pt}
    \setlength{\parskip}{6pt plus 2pt minus 1pt}}
}{% if KOMA class
  \KOMAoptions{parskip=half}}
\makeatother
\usepackage{longtable,booktabs,array}
\usepackage{calc} % for calculating minipage widths
% Correct order of tables after \paragraph or \subparagraph
\usepackage{etoolbox}
\makeatletter
\patchcmd\longtable{\par}{\if@noskipsec\mbox{}\fi\par}{}{}
\makeatother
% Allow footnotes in longtable head/foot
\IfFileExists{footnotehyper.sty}{\usepackage{footnotehyper}}{\usepackage{footnote}}
\makesavenoteenv{longtable}
\setlength{\emergencystretch}{3em} % prevent overfull lines
\providecommand{\tightlist}{%
  \setlength{\itemsep}{0pt}\setlength{\parskip}{0pt}}
\usepackage{bookmark}
\IfFileExists{xurl.sty}{\usepackage{xurl}}{} % add URL line breaks if available
\urlstyle{same}
\hypersetup{
  hidelinks,
  pdfcreator={LaTeX via pandoc}}

\author{SCX}
\date{2026}

\begin{document}

\section{Completeness Analysis of SCX
Self-Evolution}\label{completeness-analysis-of-scx-self-evolution}

\begin{quote}
\textbf{Part of the SCX Self-Evolution Theory Series} \textbf{Status}:
Formal analysis \textbar{} \textbf{Audit}: Pre-review
\textbf{Prerequisites}: THEOREMS\_UNIFIED.md (Thm 1-5, Prop 6),
self-evolution definitions (Files 1-6)
\end{quote}

\begin{center}\rule{0.5\linewidth}{0.5pt}\end{center}

\subsection{Table of Contents}\label{table-of-contents}

\begin{enumerate}
\def\labelenumi{\arabic{enumi}.}
\tightlist
\item
  \hyperref[1-finite-structure-space-argument]{Finite Structure Space
  Argument}
\item
  \hyperref[2-covering-number-analysis]{Covering Number Analysis}
\item
  \hyperref[3-effective-state-space-size-and-memory-capacity]{Effective
  State Space Size and Memory Capacity}
\item
  \hyperref[4-existence-of-finite-termination-time]{Existence of Finite
  Termination Time}
\item
  \hyperref[5-connection-to-the-physical-church-turing-thesis]{Connection
  to the Physical Church-Turing Thesis}
\item
  \hyperref[6-godel-incompleteness-analogy]{Godel Incompleteness
  Analogy}
\item
  \hyperref[7-theorem-se-2-completeness-bound]{Theorem SE-2:
  Completeness Bound}
\item
  \hyperref[8-implications-for-the-self-evolution-loop]{Implications for
  the Self-Evolution Loop}
\item
  \hyperref[9-summary]{Summary}
\end{enumerate}

\begin{center}\rule{0.5\linewidth}{0.5pt}\end{center}

\subsection{1. Finite Structure Space
Argument}\label{finite-structure-space-argument}

\subsubsection{1.1 Physical Constraints on the SCX
System}\label{physical-constraints-on-the-scx-system}

The SCX self-evolution framework operates under three universal physical
constraints that render the configuration space finite.

\textbf{Constraint 1: Finite Data.} The total volume of NEP data
available to the system is finite:
\[|\mathcal{D}_{\text{total}}| \leq N_{\text{max}} < \infty\] where
\(\mathcal{D}_{\text{total}} = \bigcup_{t=0}^{\infty} \mathcal{D}_t\) is
the union of all data observed across all evolution steps, and
\(N_{\text{max}}\) is bounded by physical storage capacity, experimental
budget, or the finite supply of NEP-relevant configurations.

\textbf{Constraint 2: Finite Compute.} Per iteration, the system expends
at most \(C_{\text{max}}\) FLOPs. Since each compute budget is finite,
only finitely many candidate gatekeeper updates \(S_t \to S_{t+1}\) can
be evaluated between successive memory-bank modifications.

\textbf{Constraint 3: Finite Precision.} All numerical quantities in the
SCX system are stored with finite machine precision
\(\varepsilon_{\text{mach}} > 0\). Two configurations whose parameters
differ by less than \(\varepsilon_{\text{mach}}\) are physically
indistinguishable.

\subsubsection{1.2 Distinguishable
Configurations}\label{distinguishable-configurations}

Let \(\mathcal{F}\) be the class of all possible gatekeeper scoring
functions \(S_t: \mathcal{X} \to [0,1]\) realizable within the SCX
framework. Under the finite-precision constraint, the set of
\textbf{distinguishable} gatekeeper functions is:

\[\mathcal{F}_{\text{dist}} = \left\{ S \in \mathcal{F} \;:\; S(x) \in \varepsilon_{\text{mach}} \cdot \mathbb{Z} \cap [0,1],\; \forall x \right\}\]

Since \(\mathcal{X}\) itself is finite under physical storage
constraints (at most \(N_{\text{max}}\) distinct inputs can be stored),
we have:

\[|\mathcal{F}_{\text{dist}}| \leq \left( \left\lfloor \frac{1}{\varepsilon_{\text{mach}}} \right\rfloor + 1 \right)^{N_{\text{max}}} < \infty\]

Similarly, the memory bank \(M_t\) is a subset of the finite set
\(\mathcal{D}_{\text{total}}\), hence:

\[|\{M_t : t \geq 0\}| \leq 2^{N_{\text{max}}} < \infty\]

\textbf{Proposition SE-3 (Finite Configuration Space).} Under physical
constraints (finite data, finite compute, finite precision), the set of
distinguishable states of the SCX self-evolution system is finite.

\emph{Proof.} The system state is fully described by the triple
\((S_t, M_t, f_{\theta_t})\), where: -
\(S_t \in \mathcal{F}_{\text{dist}}\) is the gatekeeper (finite set by
above argument) - \(M_t \subseteq \mathcal{D}_{\text{total}}\) is the
memory bank (finite set of cardinality \(\leq 2^{N_{\text{max}}}\)) -
\(f_{\theta_t}\) is the NEP student model. Under finite precision, the
parameter space \(\Theta\) is finite:
\(|\Theta| \leq (1/\varepsilon_{\text{mach}})^{\dim(\Theta)}\)

Therefore the Cartesian product is finite:
\[|\mathcal{Q}| = |\mathcal{F}_{\text{dist}}| \times 2^{N_{\text{max}}} \times |\Theta| < \infty\]
where \(\mathcal{Q}\) is the state space of the self-evolution dynamical
system. \(\square\)

\begin{center}\rule{0.5\linewidth}{0.5pt}\end{center}

\subsection{2. Covering Number Analysis}\label{covering-number-analysis}

\subsubsection{2.1 Covering Number of the Gatekeeper
Class}\label{covering-number-of-the-gatekeeper-class}

Define the gatekeeper function class \(\mathcal{F}\) as the set of all
scoring functions realizable by the SCX framework. Let \(\mathcal{F}\)
be parameterized by \(d\)-dimensional parameter vector
\(w \in \mathbb{R}^d\) (e.g., state prototypes, threshold parameters,
and weighting coefficients).

\textbf{Definition (Covering Number).} The covering number
\(N(\varepsilon, \mathcal{F}, \|\cdot\|_{\infty})\) is the minimum
number of \(\varepsilon\)-balls (in \(\|\cdot\|_{\infty}\) norm) needed
to cover \(\mathcal{F}\).

\textbf{Proposition SE-4 (Covering Number Bound).} For the SCX
gatekeeper function class \(\mathcal{F}\) with \(d\)-dimensional
parameter space contained in \([-R, R]^d\):

\[N(\varepsilon, \mathcal{F}, \|\cdot\|_{\infty}) \leq \left( \frac{4R}{\varepsilon} \right)^d \cdot (1 + o(1))\]

\emph{Proof sketch.} Following standard results from empirical process
theory (van der Vaart \& Wellner, 1996), the covering number of a
Lipschitz-parameterized function class is bounded by the covering number
of the parameter space. If \(w \mapsto S_w\) is \(L\)-Lipschitz with
respect to \(\|\cdot\|_{\infty}\) and the parameter space
\(W \subseteq \mathbb{R}^d\) is bounded, then:

\[N(\varepsilon, \mathcal{F}, \|\cdot\|_{\infty}) \leq N(\varepsilon/L, W, \|\cdot\|_2)\]

For \(W = [-R, R]^d\), the Euclidean covering number satisfies
\(N(\delta, W, \|\cdot\|_2) \leq (2R\sqrt{d}/\delta)^d\). Setting
\(\delta = \varepsilon/L\) yields:

\[N(\varepsilon, \mathcal{F}, \|\cdot\|_{\infty}) \leq \left( \frac{2RL\sqrt{d}}{\varepsilon} \right)^d\]

Absorbing constants into \(R\) gives the stated bound. \(\square\)

\subsubsection{2.2 Metric Entropy}\label{metric-entropy}

The \textbf{metric entropy} is the logarithm of the covering number:

\[\mathcal{H}(\varepsilon, \mathcal{F}) = \log N(\varepsilon, \mathcal{F}, \|\cdot\|_{\infty}) \leq d \cdot \log\left( \frac{4R}{\varepsilon} \right)\]

\textbf{Proposition SE-5 (Metric Entropy Growth).} For the SCX
gatekeeper class \(\mathcal{F}\) with fixed parameter dimension \(d\):

\[\mathcal{H}(\varepsilon, \mathcal{F}) = O\left( d \cdot \log\frac{1}{\varepsilon} \right) = \tilde{O}(d)\]

The metric entropy scales \textbf{polylogarithmically} in
\(1/\varepsilon\), not polynomially as \(O(1/\varepsilon^d)\). This is
because the gatekeeper class is \textbf{parametric}
(finite-dimensional), not nonparametric.

\textbf{Correction}: The problem statement mentions
\(O(1/\varepsilon^d)\) metric entropy. This would correspond to a
\textbf{nonparametric} class (e.g., a Sobolev or Holder class). For the
parametric SCX gatekeeper, the correct scaling is
\(O(d \log(1/\varepsilon))\). We include this correction for
mathematical accuracy. The physical implication is the same: the metric
entropy is finite for any \(\varepsilon > 0\).

\begin{center}\rule{0.5\linewidth}{0.5pt}\end{center}

\subsection{3. Effective State Space Size and Memory
Capacity}\label{effective-state-space-size-and-memory-capacity}

\subsubsection{3.1 Effective State Space
Size}\label{effective-state-space-size}

Under machine precision \(\varepsilon_{\text{mach}}\) and finite feature
dimension \(d_\phi\), the effective number of distinguishable input
points is:

\[|\mathcal{X}_{\text{eff}}| \leq \left( \frac{1}{\varepsilon_{\text{mach}}} \right)^{d_\phi}\]

This follows from discretizing each of the \(d_\phi\) feature dimensions
into at most \(1/\varepsilon_{\text{mach}}\) bins. In practice,
\(|\mathcal{X}_{\text{eff}}|\) is further bounded by \(N_{\text{max}}\)
(the total data volume).

\subsubsection{3.2 Memory Bank Capacity}\label{memory-bank-capacity}

The memory bank \(M_t\) stores triples \((x_i, y_i, S_t(x_i))\) for
select data points. Its maximum size is:

\[|M_\infty| \leq N_{\text{max}} < \infty\]

where \(M_\infty = \lim_{t\to\infty} M_t\) (the limit exists as \(M_t\)
is monotonic non-decreasing and bounded). The memory bank capacity is
\textbf{finite by construction} in the SCX framework: the total amount
of NEP data ever generated is finite.

\begin{center}\rule{0.5\linewidth}{0.5pt}\end{center}

\subsection{4. Existence of Finite Termination
Time}\label{existence-of-finite-termination-time}

\subsubsection{4.1 Monotonic Structure}\label{monotonic-structure}

The self-evolution loop has two monotonic components:

\begin{enumerate}
\def\labelenumi{\arabic{enumi}.}
\tightlist
\item
  \textbf{Memory bank monotonicity}: \(M_t \subseteq M_{t+1}\) for all
  \(t\) (data is only added, never removed).
\item
  \textbf{Gatekeeper improvement}: The sequence \(S_t\) converges to a
  fixed point under the Lyapunov function
  \(\Phi(S_t, M_t, f_{\theta_t})\) (Theorem SE-1).
\end{enumerate}

Since the set of distinguishable system states \(\mathcal{Q}\) is
finite, and the evolution is deterministic (or Markovian), any sequence
\((S_t, M_t, f_{\theta_t})\) must eventually visit a previously visited
state.

\subsubsection{4.2 Finite-Time Fixed
Point}\label{finite-time-fixed-point}

\textbf{Lemma SE-2 (Cycle Detection).} In a finite-state deterministic
dynamical system, any infinite trajectory must eventually enter a cycle.
If the dynamics are strictly improving (i.e., \(\Phi\) strictly
decreases at each step until a fixed point), the only possible cycles
are fixed points.

\emph{Proof.} This is a standard result from the theory of finite-state
machines (Moore, 1956). A deterministic map on a finite set of
cardinality \(Q\) can have at most \(Q\) distinct states before
repeating. Once a state repeats, the trajectory is periodic. If a strict
Lyapunov function exists, the value of \(\Phi\) must strictly decrease
at each step until a fixed point is reached, ruling out cycles of length
\(>1\). \(\square\)

\textbf{Theorem SE-2 (Completeness Bound).} Under the finite physical
constraints of Section 1, there exists a finite time \(T^* < \infty\)
such that for all \(t \geq T^*\):

\[\Phi(S_t, M_t, f_{\theta_t}) = \Phi(S_{T^*}, M_{T^*}, f_{\theta_{T^*}})\]

i.e., the self-evolution reaches an exact fixed point. Moreover, for any
\(\varepsilon > 0\), there exists \(T_\varepsilon^* < \infty\) such that
for all \(t \geq T_\varepsilon^*\):

\[\|\Phi(S_t, M_t, f_{\theta_t}) - \Phi_{\text{opt}}\|_\infty \leq \varepsilon\]

where \(\Phi_{\text{opt}}\) is the optimal value achievable within the
finite configuration space.

\emph{Proof.}

\textbf{Part 1 (exact fixed point).} By Proposition SE-3, the
configuration space \(\mathcal{Q}\) is finite with
\(|\mathcal{Q}| = Q < \infty\). The self-evolution update is a
deterministic mapping \(G: \mathcal{Q} \to \mathcal{Q}\). Consider the
orbit \(\{q_0, q_1, q_2, \ldots\}\) where
\(q_t = (S_t, M_t, f_{\theta_t})\). Since \(|\mathcal{Q}| = Q\), by the
pigeonhole principle, among the first \(Q+1\) states, at least two are
identical: \(q_a = q_b\) for some \(0 \leq a < b \leq Q\). If \(a = 0\),
then the system is already periodic. If \(a > 0\), then
\(q_0, \ldots, q_{a-1}\) are transient, and \(q_a, \ldots, q_{b-1}\)
form a cycle.

Now, by the strict Lyapunov property (Theorem SE-1, assumption),
\(\Phi(q_{t+1}) < \Phi(q_t)\) unless \(q_{t+1} = q_t\). In a cycle of
length \(L > 1\), we would require
\(\Phi(q_a) < \Phi(q_{a+1}) < \cdots < \Phi(q_{a+L-1}) < \Phi(q_a)\), a
contradiction. Therefore \(L = 1\), and \(q_a\) is a fixed point.
Setting \(T^* = a\) completes the proof.

\textbf{Part 2 (\(\varepsilon\)-approximate fixed point).} Even if the
strict Lyapunov condition holds only approximately (i.e., the decrease
is at most \(\delta_t\) per step with
\(\sum_{t=0}^\infty \delta_t \leq D < \infty\)), the convergence to an
\(\varepsilon\)-neighborhood of the optimal value follows from the fact
that \(\Phi\) takes values in a compact subset of \(\mathbb{R}\) (by
boundedness of the loss). The finite covering number implies that after
at most \(N(\varepsilon, \mathcal{Q}, \|\cdot\|)\) steps, the system
must be within \(\varepsilon\) of some previously visited state, and the
optimality gap must be bounded by \(\varepsilon\) due to the covering
radius. \(\square\)

\subsubsection{\texorpdfstring{4.3 Explicit Bound on
\(T^*\)}{4.3 Explicit Bound on T\^{}*}}\label{explicit-bound-on-t}

A crude but explicit bound on the termination time follows from the size
of the configuration space:

\[T^* \leq |\mathcal{Q}| = |\mathcal{F}_{\text{dist}}| \times 2^{N_{\text{max}}} \times |\Theta|\]

Using the bounds from Sections 1-3:

\[T^* \leq \left( \frac{1}{\varepsilon_{\text{mach}}} + 1 \right)^{N_{\text{max}}} \times 2^{N_{\text{max}}} \times \left( \frac{1}{\varepsilon_{\text{mach}}} \right)^{\dim(\Theta)}\]

This bound is astronomically large for realistic \(N_{\text{max}}\)
(e.g., \(10^6\)), but it remains \textbf{finite}. In practice, the
convergence is much faster due to the Lyapunov structure; the bound here
is a worst-case theoretical guarantee.

\begin{center}\rule{0.5\linewidth}{0.5pt}\end{center}

\subsection{5. Connection to the Physical Church-Turing
Thesis}\label{connection-to-the-physical-church-turing-thesis}

\subsubsection{5.1 The SCX System as a Finite-State
Machine}\label{the-scx-system-as-a-finite-state-machine}

The physical Church-Turing thesis (Deutsch, 1985) states that every
physically realizable system can be simulated by a universal Turing
machine. A corollary is that any finite physical system with finite
memory is equivalent to a \textbf{finite-state machine (FSM)}.

The SCX self-evolution system, operating on a finite computer with
finite storage, finite data, and finite precision, is precisely such a
finite physical system. Therefore:

\begin{itemize}
\tightlist
\item
  The SCX self-evolution is equivalent to a \textbf{deterministic finite
  automaton (DFA)} with at most \(|\mathcal{Q}|\) states.
\item
  The evolution loop is a state transition function
  \(\delta: \mathcal{Q} \times \mathcal{I} \to \mathcal{Q}\) where
  \(\mathcal{I}\) represents the finite set of possible NEP data inputs.
\item
  The termination guarantee (Theorem SE-2) is a direct consequence of
  the finiteness of the state machine: every DFA on a finite input
  string eventually stabilizes or enters a cycle.
\end{itemize}

\subsubsection{5.2 Implications}\label{implications}

\begin{longtable}[]{@{}
  >{\raggedright\arraybackslash}p{(\linewidth - 2\tabcolsep) * \real{0.4074}}
  >{\raggedright\arraybackslash}p{(\linewidth - 2\tabcolsep) * \real{0.5926}}@{}}
\toprule\noalign{}
\begin{minipage}[b]{\linewidth}\raggedright
Principle
\end{minipage} & \begin{minipage}[b]{\linewidth}\raggedright
SCX Implication
\end{minipage} \\
\midrule\noalign{}
\endhead
\bottomrule\noalign{}
\endlastfoot
Finite physical resources & Finite configuration space
\(|\mathcal{Q}| < \infty\) \\
Physical Church-Turing thesis & SCX evolution is computable and
simulable \\
Finite-state machine equivalence & Ultimate convergence to fixed point
guaranteed \\
Halting problem & \textbf{Not applicable}: the system has finite states,
so halting is decidable \\
\end{longtable}

\textbf{Clarification}: Unlike a general Turing machine (for which
halting is undecidable), the SCX self-evolution system operates in a
\textbf{finite} configuration space. Therefore, the question ``does the
evolution terminate?'' is trivially decidable --- it must either reach a
fixed point or cycle. The structural analysis in Theorem SE-2 shows only
fixed points are possible under the Lyapunov condition.

\begin{center}\rule{0.5\linewidth}{0.5pt}\end{center}

\subsection{6. Godel Incompleteness
Analogy}\label{godel-incompleteness-analogy}

\subsubsection{6.1 Structural Parallel}\label{structural-parallel}

Godel's first incompleteness theorem states that any sufficiently
powerful formal system contains true statements that cannot be proven
within the system. The SCX self-evolution system exhibits a structural
parallel, though it is \textbf{not} a formal instantiation of Godel's
theorem.

\textbf{The SCX system \(\mathcal{T}\) consists of:} - A data universe
\(\mathcal{D}_{\text{total}}\) (the ``axioms'') - A gatekeeper scoring
function \(S_t\) (the ``inference engine'') - A memory bank \(M_t\) (the
``theorem database'') - A NEP student \(f_{\theta_t}\) (the ``model
builder'') - Update rules \(S_t \to S_{t+1}\) and \(M_t \to M_{t+1}\)
(the ``rules of inference'')

\textbf{Claim (SE-C1): Within \(\mathcal{T}\), there exist true
statements about data quality that cannot be proven within
\(\mathcal{T}\).}

\emph{Reasoning.} The truth of a data-quality statement (e.g., ``sample
\(x\) is label noise'') is determined by the \textbf{true physical
oracle} \(f^*\), which is external to \(\mathcal{T}\). The gatekeeper
\(S_t\) provides a \textbf{provable} score (in the sense of being
computed by \(\mathcal{T}\)'s rules), but this score may differ from the
oracle's judgment. By Theorem 3 (Noise-Difficulty Unidentifiability),
there exist configurations where \(\mathcal{T}\) cannot distinguish
noise from difficulty given its finite evidence. The oracle \(f^*\)
``knows'' the truth, but \(\mathcal{T}\) cannot prove it.

\textbf{Formal parallel:}
\[\text{Godel: } \mathcal{P} \nvdash \varphi \text{ but } \mathbb{N} \models \varphi \qquad \Longleftrightarrow \qquad \text{SCX: } \mathcal{T} \nvdash \text{noise}(x) \text{ but } f^* \models \text{noise}(x)\]

where \(\mathcal{P}\) is a formal system, \(\mathbb{N}\) is the standard
model of arithmetic, \(\varphi\) is a Godel sentence, ``noise\((x)\)''
is the proposition ``sample \(x\) is label noise'', and \(\models\)
denotes truth in the intended model.

\subsubsection{6.2 Self-Reference: The System Cannot Certify Its Own
Completeness}\label{self-reference-the-system-cannot-certify-its-own-completeness}

\textbf{Claim (SE-C2): The SCX self-evolution cannot certify its own
completeness.}

\emph{Reasoning.} Suppose \(\mathcal{T}\) attempts to prove that it has
reached a complete state, i.e., that its fixed point \(S_{T^*}\) is
optimal. This would require \(\mathcal{T}\) to verify:

\[\forall x \in \mathcal{X}, \; |S_{T^*}(x) - \mathbb{P}(x \text{ is noise} \mid \text{all evidence})| \leq \varepsilon\]

But the ground truth \(f^*\) is unobserved (by definition). The only way
\(\mathcal{T}\) can assess its own completeness is by its internal
consistency metric \(\Phi\), which is a function of \(\mathcal{T}\)'s
own states. This is analogous to a formal system trying to prove its own
consistency: by Godel's second incompleteness theorem, a consistent
system cannot prove its own consistency.

The SCX parallel is not a formal theorem --- the system is not a formal
logic with Peano arithmetic --- but the \textbf{epistemic limitation} is
real: without access to an external oracle, the system cannot
distinguish between ``I have converged to the truth'' and ``I have
converged to a self-consistent but incorrect fixed point.''

\subsubsection{6.3 External Validation as Meta-System
Escape}\label{external-validation-as-meta-system-escape}

The NEP physical experiment (direct DFT or experimental validation of
selected configurations) provides the \textbf{meta-system} that escapes
the incompleteness:

\[\mathcal{T}_{\text{SCX}} \xrightarrow{\text{physical validation}} \mathcal{T}^+_{\text{SCX}}\]

where \(\mathcal{T}^+\) incorporates feedback from the physical oracle
\(f^*\). This is analogous to moving from a formal system
\(\mathcal{P}\) to a stronger system \(\mathcal{P}'\) (e.g.,
\(\text{PA} \to \text{ZFC} \to \text{ZFC} + \text{inaccessible cardinal}\))
that can prove statements the original could not.

\textbf{Practical implication}: Periodic NEP validation (running new DFT
calculations on samples selected by the gatekeeper) is not merely an
engineering consideration --- it is \textbf{epistemically necessary} for
the SCX system to escape its own completeness limitations.

\subsubsection{6.4 Two Distinct
Limitations}\label{two-distinct-limitations}

It is important to distinguish two separate completeness limitations:

\textbf{(a) Godel-like self-reference limitation (epistemic):} The
system cannot certify its own optimality without external reference.
This is a limitation of self-referential evaluation --- any system that
measures its own performance using internally generated criteria cannot
guarantee that those criteria correspond to ground truth.

\textbf{(b) Physical completeness bound (computational):} Under finite
physical resources, the system must terminate after finitely many steps
(Theorem SE-2). This is a \textbf{guarantee} of finite-time convergence,
not a limitation. The limitation is that the converged state may not be
globally optimal --- only optimal within the finite configuration space
reachable from the initial conditions.

\begin{longtable}[]{@{}
  >{\raggedright\arraybackslash}p{(\linewidth - 4\tabcolsep) * \real{0.2564}}
  >{\raggedright\arraybackslash}p{(\linewidth - 4\tabcolsep) * \real{0.3846}}
  >{\raggedright\arraybackslash}p{(\linewidth - 4\tabcolsep) * \real{0.3590}}@{}}
\toprule\noalign{}
\begin{minipage}[b]{\linewidth}\raggedright
Property
\end{minipage} & \begin{minipage}[b]{\linewidth}\raggedright
Godel-like (a)
\end{minipage} & \begin{minipage}[b]{\linewidth}\raggedright
Physical (b)
\end{minipage} \\
\midrule\noalign{}
\endhead
\bottomrule\noalign{}
\endlastfoot
Nature & Epistemic & Computational \\
Implication & Cannot self-certify & Must halt in finite time \\
Resolution & External validation (NEP experiment) & Accept local
optimality \\
Mathematical basis & Theorem 3 (unidentifiability) & Theorem SE-2
(finite configuration space) \\
Severity & Fundamental & Practical (but addressable) \\
\end{longtable}

\begin{center}\rule{0.5\linewidth}{0.5pt}\end{center}

\subsection{7. Theorem SE-2: Completeness
Bound}\label{theorem-se-2-completeness-bound}

\subsubsection{7.1 Formal Statement}\label{formal-statement}

\textbf{Theorem SE-2 (Completeness Bound).} Let the SCX self-evolution
system operate under the following conditions:

\begin{enumerate}
\def\labelenumi{\arabic{enumi}.}
\tightlist
\item
  \textbf{Finite data}:
  \(|\mathcal{D}_{\text{total}}| \leq N_{\text{max}} < \infty\)
\item
  \textbf{Finite precision}: All numerical quantities are stored with
  precision \(\varepsilon_{\text{mach}} > 0\)
\item
  \textbf{Finite parameterization}: The gatekeeper \(S_t\) and NEP
  student \(f_{\theta_t}\) are parameterized by at most \(d\) real
  parameters
\item
  \textbf{Lyapunov descent}: The Lyapunov function
  \(\Phi(S_t, M_t, f_{\theta_t})\) satisfies
  \(\Phi(q_{t+1}) \leq \Phi(q_t) - \gamma_t\) for all \(t\), where
  \(\gamma_t \geq 0\) and \(\gamma_t > 0\) whenever \(q_{t+1} \neq q_t\)
\item
  \textbf{Bounded Lyapunov}: \(0 \leq \Phi(\cdot) \leq \Phi_0 < \infty\)
\end{enumerate}

Then:

\textbf{(a) Exact fixed point.} There exists \(T^* < \infty\) such that
\(q_t = q_{T^*}\) for all \(t \geq T^*\), i.e., the system reaches an
exact fixed point in finite time.

\textbf{(b) \(\varepsilon\)-approximate fixed point.} For any
\(\varepsilon > 0\), there exists \(T_\varepsilon^* < \infty\) such that
for all \(t \geq T_\varepsilon^*\):

\[\Phi(q_t) \leq \Phi_{\text{opt}} + \varepsilon\]

where \(\Phi_{\text{opt}} = \inf_{q \in \mathcal{Q}} \Phi(q)\) is the
minimal achievable Lyapunov value.

\textbf{(c) Termination bound.} The worst-case termination time
satisfies:

\[T^* \leq \left\lceil \frac{\Phi_0}{\gamma_{\min}} \right\rceil \leq \left\lceil \Phi_0 \cdot \frac{|\mathcal{Q}|}{\Phi_0 - \Phi_{\text{opt}}} \right\rceil\]

where \(\gamma_{\min} = \min\{\gamma_t : \gamma_t > 0\}\) is the minimum
positive descent step.

\emph{Proof.}

\textbf{(a)} By Proposition SE-3, \(|\mathcal{Q}| = Q < \infty\).
Consider the sequence \(\{q_0, q_1, \ldots, q_Q\}\). If all \(Q+1\)
states are distinct, then by the pigeonhole principle applied to \(Q\)
possible values of \(\Phi\) (on a finite grid induced by
\(\varepsilon_{\text{mach}}\)), at least two must have equal \(\Phi\)
values, violating strict descent. More formally:

Since \(\Phi\) can take at most
\(M_\Phi = \lceil \Phi_0/\varepsilon_{\text{mach}} \rceil\) distinct
values (under finite precision), and each step with \(q_{t+1} \neq q_t\)
strictly decreases \(\Phi\) by at least \(\varepsilon_{\text{mach}}\)
(the smallest distinguishable change), the number of non-fixed-point
steps is bounded by \(M_\Phi\). Therefore:

\[T^* \leq M_\Phi \leq \left\lceil \frac{\Phi_0}{\varepsilon_{\text{mach}}} \right\rceil < \infty\]

\textbf{(b)} Let \(\Phi_{\text{opt}}\) be the global minimum of \(\Phi\)
over \(\mathcal{Q}\). For any \(\varepsilon > 0\), define
\(\mathcal{Q}_\varepsilon = \{q \in \mathcal{Q} : \Phi(q) - \Phi_{\text{opt}} > \varepsilon\}\).
Since \(\mathcal{Q}\) is finite, \(|\mathcal{Q}_\varepsilon|\) is
finite. The system can visit at most \(|\mathcal{Q}_\varepsilon|\)
states in \(\mathcal{Q}_\varepsilon\) before either: (i) entering the
\(\varepsilon\)-optimal set
\(\mathcal{Q} \setminus \mathcal{Q}_\varepsilon\), or (ii) cycling.
Strict descent rules out cycling with
\(\Phi > \Phi_{\text{opt}} + \varepsilon\). Therefore
\(T_\varepsilon^* \leq |\mathcal{Q}_\varepsilon| \leq |\mathcal{Q}|\).

\textbf{(c)} The descent condition gives:

\[\Phi(q_0) - \Phi(q_t) = \sum_{k=0}^{t-1} \gamma_k \geq t \cdot \gamma_{\min}\]

where \(\gamma_{\min} = \min\{\gamma_t > 0\}\). Since
\(\Phi(q_t) \geq \Phi_{\text{opt}}\), we have
\(\Phi(q_0) - \Phi_{\text{opt}} \geq T^* \cdot \gamma_{\min}\), hence:

\[T^* \leq \frac{\Phi_0 - \Phi_{\text{opt}}}{\gamma_{\min}} \leq \left\lceil \frac{\Phi_0}{\gamma_{\min}} \right\rceil\]

The second inequality follows from
\(\gamma_{\min} \geq \varepsilon_{\text{mach}}\) (any distinguishable
descent is at least machine precision). The rightmost bound uses
\(|\mathcal{Q}|\) as a coarse overestimate. \(\square\)

\subsubsection{7.2 Proof Sketch Using Covering Number + Monotonic Memory
+ Bounded
Improvement}\label{proof-sketch-using-covering-number-monotonic-memory-bounded-improvement}

An alternative proof emphasizing the covering-number perspective:

\textbf{Step 1}: The covering number
\(N(\varepsilon, \mathcal{F}, \|\cdot\|_\infty) \leq (4R/\varepsilon)^d\)
implies that the gatekeeper function class can only achieve
\(N(\varepsilon, \mathcal{F}, \|\cdot\|_\infty)\) distinct
\(\varepsilon\)-separated functions.

\textbf{Step 2}: The memory bank \(M_t\) is monotonic
(\(M_t \subseteq M_{t+1}\)) and bounded (\(|M_t| \leq N_{\text{max}}\)).
Therefore it can change at most \(N_{\text{max}}\) times (it can only
grow by adding data points, one at a time).

\textbf{Step 3}: The NEP student \(f_{\theta_t}\) is updated only when
\(M_t\) changes or the gatekeeper changes. Hence the total number of
student updates is bounded by
\(|\mathcal{F}_{\text{dist}}| + N_{\text{max}}\).

\textbf{Step 4}: Each Lyapunov descent step reduces \(\Phi\) by at least
the minimum distinguishable amount \(\varepsilon_{\text{mach}}\). Since
\(\Phi\) is bounded below by \(0\), the number of descent steps is
bounded by \(\Phi_0/\varepsilon_{\text{mach}}\).

\textbf{Step 5}: Combining steps 1-4, the total number of evolution
steps before reaching a fixed point is at most:

\[T^* \leq \min\left( \frac{\Phi_0}{\varepsilon_{\text{mach}}}, \; N_{\text{max}} \cdot \left( \frac{4R}{\varepsilon_{\text{mach}}} \right)^d \right)\]

This is finite because all terms on the right-hand side are finite.

\subsubsection{7.3 Tightness Discussion}\label{tightness-discussion}

The bound \(T^* \leq \Phi_0/\varepsilon_{\text{mach}}\) is tight up to
constant factors: if each step achieves exactly the minimum descent
\(\varepsilon_{\text{mach}}\), the system requires exactly
\(\Phi_0/\varepsilon_{\text{mach}}\) steps to converge from the maximal
to the minimal Lyapunov value. The covering-number bound
\(N_{\text{max}} \cdot (4R/\varepsilon_{\text{mach}})^d\) is looser but
captures the dependence on the complexity of the gatekeeper class.

\begin{center}\rule{0.5\linewidth}{0.5pt}\end{center}

\subsection{8. Implications for the Self-Evolution
Loop}\label{implications-for-the-self-evolution-loop}

\subsubsection{8.1 Practical Consequences}\label{practical-consequences}

\begin{longtable}[]{@{}
  >{\raggedright\arraybackslash}p{(\linewidth - 4\tabcolsep) * \real{0.2600}}
  >{\raggedright\arraybackslash}p{(\linewidth - 4\tabcolsep) * \real{0.3600}}
  >{\raggedright\arraybackslash}p{(\linewidth - 4\tabcolsep) * \real{0.3800}}@{}}
\toprule\noalign{}
\begin{minipage}[b]{\linewidth}\raggedright
Implication
\end{minipage} & \begin{minipage}[b]{\linewidth}\raggedright
Theoretical Basis
\end{minipage} & \begin{minipage}[b]{\linewidth}\raggedright
Practical Meaning
\end{minipage} \\
\midrule\noalign{}
\endhead
\bottomrule\noalign{}
\endlastfoot
\textbf{Termination guarantee} & Theorem SE-2 & The self-evolution loop
will not run indefinitely; a stopping condition exists \\
\textbf{Fixed point} & Finite configuration space & The gatekeeper,
memory, and student stabilize; no further improvement possible without
external input \\
\textbf{Non-uniqueness} & Finite configuration space may have multiple
local optima & Different initial conditions may lead to different fixed
points \\
\textbf{External reset necessary} & Godel-like limitation & To escape a
suboptimal fixed point, new external data (NEP calculations) must be
introduced \\
\textbf{No infinite regress} & Physical Church-Turing thesis & The
system cannot self-improve arbitrarily many times; there is a finite
``improvement capacity'' \\
\end{longtable}

\subsubsection{8.2 What Theorem SE-2 Does NOT
Guarantee}\label{what-theorem-se-2-does-not-guarantee}

\begin{enumerate}
\def\labelenumi{\arabic{enumi}.}
\tightlist
\item
  \textbf{Global optimality}: The fixed point is optimal only within the
  reachable configuration space, not necessarily globally optimal.
\item
  \textbf{Correctness}: The fixed point may be incorrect if the
  gatekeeper converges to a scoring function that mislabels noise as
  clean (or vice versa). Theorem 3 shows this is unavoidable in
  principle.
\item
  \textbf{Uniqueness}: The fixed point depends on initial conditions
  (initial gatekeeper, initial memory, initial student). Different
  trajectories may converge to different fixed points.
\item
  \textbf{Speed}: The bound
  \(T^* \leq \Phi_0/\varepsilon_{\text{mach}}\) is astronomically loose
  for practical \(\varepsilon_{\text{mach}}\). The actual convergence
  time depends on the Lyapunov descent rate, which is problem-dependent.
\end{enumerate}

\subsubsection{8.3 Relation to Other
Theorems}\label{relation-to-other-theorems}

\begin{longtable}[]{@{}
  >{\raggedright\arraybackslash}p{(\linewidth - 2\tabcolsep) * \real{0.3000}}
  >{\raggedright\arraybackslash}p{(\linewidth - 2\tabcolsep) * \real{0.7000}}@{}}
\toprule\noalign{}
\begin{minipage}[b]{\linewidth}\raggedright
Theorem
\end{minipage} & \begin{minipage}[b]{\linewidth}\raggedright
Relationship to SE-2
\end{minipage} \\
\midrule\noalign{}
\endhead
\bottomrule\noalign{}
\endlastfoot
Thm 1 (Noise Detection) & Provides the gradient signal for the Lyapunov
descent; without detectable noise, \(\gamma_t = 0\) and evolution
stalls \\
Thm 2 (Weak Feature) & Bounds the quality of the fixed point: if
features are weak, even the optimal fixed point cannot outperform the
loss baseline \\
Thm 3 (Unidentifiability) & Shows that the fixed point may converge to a
self-consistent but incorrect configuration; the system cannot detect
this internally \\
Thm 4' (Exact Constant) & Characterizes the optimal detection rate at
the fixed point, linking the Lyapunov optimum to the noise detection F1
bound \\
Thm 5 (Cluster Consistency) & Governs whether the state partition (a
component of \(S_t\) and \(\mathcal{F}\)) converges to the true
partition \\
Prop 6 (Stability Diagnostic) & Provides a practical test for whether
the fixed point is meaningfully better than the baseline \\
Theorem SE-1 (Convergence) & Establishes the Lyapunov structure that
powers SE-2's descent guarantee \\
\end{longtable}

\begin{center}\rule{0.5\linewidth}{0.5pt}\end{center}

\subsection{9. Summary}\label{summary}

The completeness analysis establishes three key results:

\begin{enumerate}
\def\labelenumi{\arabic{enumi}.}
\item
  \textbf{Finite Structure Guarantee (Proposition SE-3)}: Under physical
  constraints, the SCX self-evolution system operates in a finite
  configuration space. This is a consequence of finite data, finite
  compute, and finite precision.
\item
  \textbf{Termination Guarantee (Theorem SE-2)}: The system reaches a
  fixed point in finite time. The maximum number of steps is bounded by
  \(\Phi_0/\varepsilon_{\text{mach}}\) (minimum descent steps) or, more
  coarsely, by the cardinality of the configuration space.
\item
  \textbf{Self-Limitation of Self-Certification (Claims SE-C1, SE-C2)}:
  Despite the termination guarantee, the system cannot internally
  certify that its fixed point is optimal or correct. External NEP
  validation provides the necessary meta-system escape.
\end{enumerate}

These results together establish that the SCX self-evolution framework
is \textbf{computationally well-founded} (it terminates) but
\textbf{epistemically bounded} (it cannot verify its own success). This
mirrors similar completeness-incompleteness trade-offs in other areas of
mathematics and computer science, and it provides a principled
justification for the role of physical validation in the SCX pipeline.

\begin{center}\rule{0.5\linewidth}{0.5pt}\end{center}

\subsection{References}\label{references}

\begin{enumerate}
\def\labelenumi{\arabic{enumi}.}
\tightlist
\item
  Deutsch, D. (1985). Quantum theory, the Church-Turing principle and
  the universal quantum computer. \emph{Proceedings of the Royal Society
  of London A}, 400(1818), 97-117.
\item
  Godel, K. (1931). Uber formal unentscheidbare Satze der Principia
  Mathematica und verwandter Systeme I. \emph{Monatshefte fur Mathematik
  und Physik}, 38(1), 173-198.
\item
  Moore, E. F. (1956). Gedanken-experiments on sequential machines.
  \emph{Automata Studies}, Annals of Mathematical Studies 34, 129-153.
\item
  van der Vaart, A. W., \& Wellner, J. A. (1996). \emph{Weak Convergence
  and Empirical Processes}. Springer.
\item
  Wolfram, S. (2002). \emph{A New Kind of Science}. Wolfram Media.
  {[}Chapter 3: The Principle of Computational Equivalence{]}
\item
  Aaronson, S. (2013). \emph{Quantum Computing Since Democritus}.
  Cambridge University Press. {[}Chapter 10: Godel, Turing, and
  Friends{]}
\item
  Cover, T. M., \& Thomas, J. A. (2006). \emph{Elements of Information
  Theory} (2nd ed.). Wiley.
\item
  Turing, A. M. (1936). On computable numbers, with an application to
  the Entscheidungsproblem. \emph{Proceedings of the London Mathematical
  Society}, 2(42), 230-265.
\end{enumerate}

\begin{center}\rule{0.5\linewidth}{0.5pt}\end{center}

\emph{End of 07\_completeness.md}

\end{document}
